\section*{Chapitre \ref{ch:mx:fees-rebates-and-venue-economics} --- Frais, rebates et économie des plateformes}

\begin{solution}{exo:mx:fees-rebates-and-venue-economics:1}
BZX~: $20.00+0.0030=20.0030$~; BYX~: $20.00-0.0002=19.9998$. La plateforme inversée fait économiser 0.32 cent par action, soit 3.20~\$ sur 1\,000 actions.
\end{solution}

\begin{solution}{exo:mx:fees-rebates-and-venue-economics:2}
Les cours vendeurs nets de frais s'entrelacent avec des écarts de $0.30-(-0.02)=0.32$ cent et $1-0.32=0.68$ cent~: le \omterm{def:mx:fees-rebates-and-venue-economics:efftick}{pas de cotation effectif} vaut 0.32 cent.
\end{solution}

\begin{solution}{exo:mx:fees-rebates-and-venue-economics:3}
INV~: $0.3\times(0.5-0.1-0.2)=0.06$ cent~; MT~: $0.1\times(0.5-0.6+0.16)=0.006$ cent.
\end{solution}

\begin{solution}{exo:mx:fees-rebates-and-venue-economics:4}
Résoudre $0.1\,(-0.1-m)=0.06$~: $m=-0.7$, soit un rebate de 0.70 cent. La plateforme verserait 0.70 au fournisseur de liquidité et recevrait 0.30 du preneur, ce qui lui ferait perdre 0.40 cent sur
chaque action échangée~: c'est impossible.
\end{solution}

\begin{solution}{exo:mx:fees-rebates-and-venue-economics:5}
Les frais totaux restent à 0.14 cent. Sans tick, le cours vendeur passe à 10.012~: le preneur paie $10.012+0.001=10.013$ comme auparavant ($10.01+0.003$), le fournisseur de liquidité touche
$10.012-0.0004=10.0116$, comme auparavant ($10.01+0.0016$). Avec un tick d'un centime, le cours vendeur reste à 10.01 (le preneur paie 10.011, le fournisseur perd 0.2 cent) ou passe à 10.02
(le preneur paie 10.021)~: le partage n'est plus \omterm{def:mx:fees-rebates-and-venue-economics:neutral}{neutre}.
\end{solution}

\begin{solution}{exo:mx:fees-rebates-and-venue-economics:6}
La capture correspond à la demi-fourchette, identique au même prix sur les deux plateformes. Ce qui change, c'est quelles transactions atteignent la file de MT~: les preneurs sensibles aux frais frappent d'abord INV,
si bien que les ordres sur MT ne sont exécutés qu'une fois la file d'INV épuisée, par de gros ordres ou des ordres informés après lesquels le prix se déplace. Plus les preneurs sont sensibles aux frais, plus les
exécutions sur MT sont de ce type.
\end{solution}

\begin{solution}{exo:mx:fees-rebates-and-venue-economics:7}
Trois graines de deux heures~: la part exécutée de MT (8.4~\%) et la sélection adverse (0.63 cent) ne changent pas, car le rebate de prise d'INV rend toujours la prise
moins coûteuse et l'ordre de routage reste le même~; la valeur par action postée passe de 0.006 à $-0.003$ cent, soit les 0.10 cent de rebate perdus sur chaque action exécutée.
Un plafond plus bas réduit la marge pour les rebates sans modifier qui est exécuté en premier.
\end{solution}

\begin{solution}{exo:mx:fees-rebates-and-venue-economics:8}
Le rebate revient au courtier, et l'ordre se place là où les preneurs sensibles aux frais arrivent en dernier~: il est exécuté moins souvent et, lorsqu'il l'est, plus souvent face à des flux
qui continuent de faire bouger le prix. Battalio, Corwin et Jennings ont mis en évidence une relation négative entre la qualité d'exécution des \omterm{def:mx:the-limit-order-book:orders}{ordres à cours limité} et le niveau du rebate~; la
simulation donne des taux d'exécution de 8.4~\% contre 28.2~\% et une sélection adverse de 0.63 contre 0.13 cent. Il faut mesurer les taux d'exécution et les mark-outs par plateforme avant de conclure
à une baisse des coûts.
\end{solution}

\begin{solution}{pb:mx:fees-rebates-and-venue-economics:1}
\textbf{1.} BZX~: rebate de 0.16 cent pour ajouter, frais de 0.30 pour retirer~; BYX~: frais de 0.20 pour ajouter, rebate de 0.02 pour retirer.
\textbf{2.} 20.0030 et 19.9998.
\textbf{3.} Le plus petit écart entre les prix nets de frais d'une plateforme à l'autre~: 0.32 cent.
\textbf{4.} Les plateformes ne peuvent pas coter entre les ticks~; des frais différents leur permettent d'offrir des prix infra-tick nets de frais (Chao, Yao et Ye).
\textbf{5.} Voir la proposition~: $a+f^{take}=R+F$ sans tick~; avec un tick, $a$ est arrondi au niveau supérieur de la grille.
\textbf{6.} Sans tick, le cours vendeur passe à 10.012 et rien d'autre ne change~; avec un tick, il reste à 10.01 ou passe à 10.02.
\textbf{7.} Sur la position dans la file~: les fournisseurs de liquidité se disputent la place en rejoignant les files, si bien que le partage modifie la longueur des files.
\textbf{8.} Les cotations ajustées et les coûts des demandeurs, frais inclus, n'ont pas changé~; les \omterm{def:mx:decomposing-the-spread:quoted}{fourchettes affichées} se sont resserrées et les ordres agressifs sont devenus plus fréquents.
\textbf{9.} Un prix efficient avec sauts, des \omterm{def:mx:the-limit-order-book:orders}{ordres à cours limité} au meilleur cours ou derrière, annulés à taux constant et quand ils sont périmés, du bruit et des \omterm{def:mx:the-limit-order-book:orders}{ordres au marché} informés~;
ceux qui sont sensibles aux frais routent selon le prix net de frais, ceux qui y sont insensibles selon la taille affichée.
\textbf{10.} INV 77~\% du volume, 28.2~\% exécutés, 4.6 secondes~; MT 23~\%, 8.4~\%, 8.8 secondes.
\textbf{11.} 0.63 cent sur MT contre 0.13 sur INV~: les ordres sur MT sont exécutés lorsque la file d'INV est vide, c'est-à-dire quand arrivent les gros flux ou les flux informés.
\textbf{12.} 0.006 cent sur MT, 0.057 sur INV.
\textbf{13.} Un rebate de 0.77 cent, au-delà des frais de prise de 0.30 cent~: non.
\textbf{14.} La file de MT est d'abord prise par les routeurs fondés sur la taille~: 0.052 cent sur MT contre 0.028 sur INV, et les frais d'ajout indifférents sur MT sont à peu près nuls.
\textbf{15.} 0.074 contre 0.015 cent.
\textbf{16.} Rule 610T, adoptée le 19 décembre 2018~: 1\,460 actions réparties entre un groupe à plafond de frais d'accès de 0.0010~\$ et un groupe sans rebate~; annulée par la D.C. Circuit le
16 juin 2020.
\textbf{17.} Des courtiers envoyant des \omterm{def:mx:the-limit-order-book:orders}{ordres à cours limité} vers les plateformes à rebate le plus élevé, et une qualité d'exécution des \omterm{def:mx:the-limit-order-book:orders}{ordres à cours limité} plus faible là où les rebates sont plus élevés.
\textbf{18.} Les rebates doivent rester inférieurs aux frais de prise de 0.10 cent, de sorte que le montant maximal qu'une plateforme maker-taker peut verser baisse de 0.20 cent.
\textbf{19.} \emph{Résultat nommé}~: lorsque tous les preneurs sont sensibles aux frais, une action postée est exécutée dans 8.4~\% des cas sur la plateforme maker-taker, avec 0.63 cent de
sélection adverse sur 30 secondes, contre 28.2~\% et 0.13 cent sur la plateforme inversée~; un fournisseur de liquidité ne serait indifférent qu'à un rebate de 0.77 cent, que la plateforme ne peut pas
verser~; si la moitié des preneurs routent selon la taille, les frais d'ajout indifférents sont à peu près nuls.
\textbf{20.} Une place au fond de la file consolidée, dont la valeur dépend de ce que le routage des preneurs lui accorde.
\end{solution}

\begin{solution}{iq:mx:fees-rebates-and-venue-economics:1}
La plateforme maker-taker rémunère le côté passif et facture le preneur~; la plateforme inversée fait l'inverse. Pour être exécuté rapidement, il faut poster sur la plateforme inversée~: les preneurs sensibles aux frais
la frappent en premier à prix égal, et sa file est plus courte.
\iqlookfor{Rôles des frais~; ordre de routage~; longueur de la file.}
\end{solution}

\begin{solution}{iq:mx:fees-rebates-and-venue-economics:2}
Sans tick, oui~: les cotations absorbent le mouvement et les prix nets de frais ne dépendent que des frais totaux. Avec un tick, les cotations ne peuvent pas bouger d'une fraction de celui-ci, donc
le partage modifie les prix nets de frais et la longueur des files~; il compte aussi lorsque certains traders ne paient pas directement de frais (flux de particuliers, certains courtiers).
\iqlookfor{Neutralité du prix continu~; le tick~; qui supporte les frais.}
\end{solution}

\begin{solution}{iq:mx:fees-rebates-and-venue-economics:3}
Des flux différents atteignent la plateforme (ordre de routage, preneurs sensibles aux frais frappant d'abord une autre plateforme), ou des teneurs de cotations différents ou des types d'ordres différents. Il faut conditionner sur les
mêmes transactions~: comparer les exécutions issues des mêmes balayages agressifs, ou randomiser au hasard la plateforme des ordres passifs et comparer les mark-outs.
\iqlookfor{Ordre de routage~; sélection~; test randomisé.}
\end{solution}

\begin{solution}{iq:mx:fees-rebates-and-venue-economics:4}
Ces ordres deviennent le fond de la file consolidée~: ils ne sont exécutés que lorsque les plateformes moins chères sont épuisées, moins souvent et face à des flux plus informés, et
leur rebate doit compenser cela.
\iqlookfor{Ordre des files entre plateformes~; sélection adverse.}
\end{solution}

\begin{solution}{iq:mx:fees-rebates-and-venue-economics:5}
Affecter au hasard des actions à des régimes de frais avec un groupe témoin (comme le visait le Transaction Fee Pilot), mesurer les fourchettes, la profondeur, les taux d'exécution et les mark-outs,
les parts de routage et les volumes, et estimer des \omterm{def:mx:tick-size-queues-and-priority:did}{doubles différences} (difference-in-differences) avec des erreurs types groupées (clustered)~; vérifier les effets de débordement entre plateformes.
\iqlookfor{Randomisation~; témoin~; débordements.}
\end{solution}

\begin{solution}{iq:mx:fees-rebates-and-venue-economics:6}
Les rebates doivent descendre sous 0.10 cent, de sorte que l'avantage des plateformes maker-taker pour les fournisseurs de liquidité se réduit, que les files s'y raccourcissent et que l'avantage
des plateformes inversées pour les preneurs s'amenuise~; mon P\&L perd le revenu des rebates et doit provenir de la capture du spread et de la position dans la file, si bien que je recote là où les exécutions sont moins toxiques.
\iqlookfor{Arithmétique des frais~; réponse des files~; composantes du P\&L.}
\end{solution}
